\textbf{Program 3: Complete infix-to-postfix converter.}

\textbf{File:} \texttt{infix\_to\_postfix.cpp}.

\textbf{Chapter connection:} Lab-manual application: infix to postfix.

Implement the chapter token rules with a stack of operators. precedence and rightAssociative decide when to pop. The trace records the stack and output after each token. Operands are single letters/digits; this is not a full expression grammar validator.

\begin{adslabfigure}[How this program runs]
\begin{adslabgraphic}
\begin{tikzpicture}[font=\small,execute at begin node={\hyphenpenalty=10000\exhyphenpenalty=10000}]
\node[utpstep,text width=2.50cm,minimum height=1.15cm] (s0) at (0.0,0) {\textbf{1}\\Token / helpers};
\node[utpstep,text width=2.50cm,minimum height=1.15cm] (s1) at (2.95,0) {\textbf{2}\\Operands to output};
\node[utpstep,text width=2.50cm,minimum height=1.15cm] (s2) at (5.9,0) {\textbf{3}\\Operators on stack};
\node[utpstep,text width=2.50cm,minimum height=1.15cm] (s3) at (8.85,0) {\textbf{4}\\Drain / test cases};
\draw[utparrow] (s0.east)--(s1.west);
\draw[utparrow] (s1.east)--(s2.west);
\draw[utparrow] (s2.east)--(s3.west);
\end{tikzpicture}
\end{adslabgraphic}
\end{adslabfigure}


\begin{utpsteps}\setcounter{enumi}{0}\item \textbf{Headers and initial definitions.}\end{utpsteps}

\texttt{infix\_to\_postfix.cpp} --- lines 14--25.

\lstinputlisting[firstline=14,lastline=25]{all_experiments/Lab05/infix_to_postfix.cpp}


\tcbbreak

\begin{utpsteps}\setcounter{enumi}{1}\item \textbf{Implement isOperator, rightAssociative, toPostfix.}\end{utpsteps}

\texttt{infix\_to\_postfix.cpp} --- lines 26--61.

% Copied from chapters/05-stacks.tex (fig:infix-trace).
% Keep the chapter drawing, type sizes, colours, and node styles unchanged.
\begin{adslabfigure}[Infix to postfix conversion traced]
\begin{adslabgraphic}
\begin{tikzpicture}[font=\small]
\node[anchor=north west,inner sep=0pt] at (0,0){
\begin{tabular}{@{}>{\bfseries\color{UTPNavy}}c
                   >{\raggedright\arraybackslash}m{3.9cm}
                   >{\ttfamily\color{UTPBlue}}l
                   >{\ttfamily\color{UTPGreen!70!black}}l@{}}
\toprule
\rowcolor{UTPPaleBlue!92}
\utpth{Read} & \utpth{Action} & \utpth{Stack} & \utpth{Output} \\
\midrule
A & operand: emit            & \phantom{x}      & A          \\
+ & stack empty: push        & +                & A          \\
B & operand: emit            & +                & A B        \\
* & higher than +: push      & + *              & A B        \\
C & operand: emit            & + *              & A B C      \\
- & pop * then +, push -     & -                & A B C * +  \\
D & operand: emit            & -                & A B C * + D\\
\midrule
\textit{end} & pop remaining & \phantom{x}      & A B C * + D -\\
\bottomrule
\end{tabular}};
\end{tikzpicture}
\end{adslabgraphic}
\end{adslabfigure}


\lstinputlisting[firstline=26,lastline=61]{all_experiments/Lab05/infix_to_postfix.cpp}


\tcbbreak

\begin{utpsteps}\setcounter{enumi}{2}\item \textbf{Connect the test cases in main.}\end{utpsteps}

\texttt{infix\_to\_postfix.cpp} --- lines 62--97.

\lstinputlisting[firstline=62,lastline=97]{all_experiments/Lab05/infix_to_postfix.cpp}


\tcbbreak

\begin{utpsteps}\setcounter{enumi}{3}\item \textbf{Connect the test cases in main.}\end{utpsteps}

\texttt{infix\_to\_postfix.cpp} --- lines 98--109. Continue in the same file; do not add another header or opening brace.

\lstinputlisting[firstline=98,lastline=109]{all_experiments/Lab05/infix_to_postfix.cpp}


\tcbbreak

\textbf{Build, run, and inspect infix\_to\_postfix.}

From this program\textquotesingle s folder, run these commands in PowerShell:

\begin{lstlisting}[style=utpterminal,language={}]

g++ -std=c++17 -Wall -Wextra infix_to_postfix.cpp -o infix_to_postfix.exe

.\infix_to_postfix.exe

\end{lstlisting}

\textbf{Expected output excerpt.} Compare values, order, and pass/fail messages; spacing need not match.

\begin{lstlisting}[style=utpterminal,language={}]

== Trace of A+B*C-D (the Chapter 5 figure) ==
  read A   stack []   output A
  read +   stack [+]   output A
  read B   stack [+]   output AB
  read *   stack [+*]   output AB
  read C   stack [+*]   output ABC
  read -   stack [-]   output ABC*+
  read D   stack [-]   output ABC*+D
result: ABC*+D-

== Test expressions ==
  [ok]   A+B*C  ->  ABC*+   (expected ABC*+)
  [ok]   (A+B)*C  ->  AB+C*   (expected AB+C*)
  [ok]   A+B*C-D  ->  ABC*+D-   (expected ABC*+D-)
  [ok]   A*(B+C)/D  ->  ABC+*D/   (expected ABC+*D/)

\end{lstlisting}

The complete captured run is in \texttt{labs/outputs/infix\_to\_postfix.txt}.

\begin{tryit}

Compare A-B-C with A\textasciicircum{}B\textasciicircum{}C: expect AB-C- and ABC\textasciicircum{}\textasciicircum{}. Test (A+B and A+B) and check that failure is reported; do not use their incomplete outputs.

\end{tryit}
